\section{Preliminaries}
\label{sec:prelim}

\subsection{Problem: static call graphs for Python}

A \emph{call graph} of a program $P$ is a directed graph whose nodes are the
modules, functions, methods and lambdas defined in $P$, and whose edges
$A \rightarrow B$ state that the code of $A$ can call $B$ when $A$ runs. A
\emph{static} call-graph generator computes such a graph without running $P$.
For Python this is hard: callees are first-class values that flow through
variables, attributes, containers, parameters and return values; methods are
looked up dynamically through the C3 method resolution order (MRO); many calls
are \emph{implicit} (\texttt{\_\_init\_\_} on instantiation, \texttt{\_\_iter\_\_}
in a \texttt{for} loop, operators, properties, decorators); and libraries call
code back (\texttt{sorted(xs, key=f)}, template methods of a framework base
class).

We adopt the exact semantics of the benchmark that we evaluate on
(Section~\ref{sec:setup}). Nodes have dotted names (\texttt{pkg.mod.Class.method},
\texttt{main.<lambda1>}); classes, class bodies and comprehensions are not
nodes, and their calls belong to the enclosing function or module. A callee
outside the analysed package (a builtin, the standard library, a third-party
library) is named by a public path, e.g.\ \texttt{<builtin>.str.join},
\texttt{os.path.join} or \texttt{re.Pattern.match}. For real packages an edge
exists if \emph{some} client of the public API can make the call on a supported
Python version, so the ground truth is a may-call graph over all inputs, not
the trace of one run. Following the PyCG evaluation, we score edges whose
callee is inside the package (\emph{inside}, the headline) separately from
edges to outside code (\emph{outside}), whose names depend on how a tool models
libraries. We report precision $P = |T \cap G| / |T|$ and recall
$R = |T \cap G| / |G|$ of the edge set $T$ of a tool against the ground truth
$G$, and $F_1$.

\subsection{Baseline: PyCG}

PyCG~\cite{salis2021pycg} is the reference static call-graph generator for
Python and the strongest open-source baseline in the comparison of its
evaluation repository (\textsc{Jarvis}, pyan3, PyAnalyzer, HeaderGen and
CPyGraph are lower in both precision and recall). We use \emph{two} versions,
both executable and both run by the official runner of the evaluation:

\textbf{PyCG-2021} (the analysis of the ICSE 2021 paper, commit
\texttt{046479e}, ported to Python~3.12+ by its authors). It builds an
\emph{assignment graph} between program identifiers (variables, functions,
classes, modules, attributes), processes imports module by module, and
resolves a call by following the assignment relations from the called
expression to the functions and classes that may reach it, until a fixpoint.
It is flow- and context-insensitive and computes MROs for method lookup, but it
models neither the standard library nor the builtins beyond their names, so
calls through library values are mostly lost.

\textbf{PyCG-2026} (the redesigned engine, commit \texttt{c5c6a89}). In our
reading of its design document, it keeps the idea of a whole-program
inclusion-based (Andersen-style) value flow but over \emph{objects} named by
their creation site, with points-to sets for variables, attributes, container
elements and returns. Names local to a scope are followed by definition (a
limited form of flow sensitivity), statically dead code is skipped, outside
code is typed from the \emph{typeshed} stubs~\cite{typeshed} (so \texttt{re.match(...).group()}
resolves to \texttt{re.Match.group}), annotations and ``entry point''
receivers supply the values that only clients could pass, and implicit calls
of the interpreter and the special methods that builtins call are modelled.
It deliberately does not model what outside code calls back, has no calling
context and no flow-sensitive heap. Both versions read code and never import
or run it.

\textbf{How we use the baseline.} We run both engines unmodified, through
\texttt{evaluation/common/run.py} of the evaluation repository (each analysis
in a fresh interpreter, \texttt{PYTHONHASHSEED=0}, 30-minute limit per
program), on exactly the same entry points and package roots that the
benchmark prescribes, and score their graphs with the official scorer. Our
reproduction matches the published numbers of PyCG-2026 exactly (micro
P\,0.925/R\,0.936; 30 packages inside P\,0.945/R\,0.878), which validates the
pipeline. PyCG-2026 is a demanding baseline: its authors developed it against
this benchmark, so it is likely tuned to the benchmark's conventions, a point
we return to in Section~\ref{sec:threats}.
